% SPDX-License-Identifier: MIT
% Copyright (c) 2026 Evgeny Baulin
% Optimization Methods, Seminar 06: one-page cheat sheet (English).
% Generated by checks/06. Second-Order Methods Part Two.ipynb from web/06/js/content-cheatsheet.js and content/students.json. Do not edit by hand.
% Compile with: latexmk -xelatex -interaction=nonstopmode -halt-on-error Cheatsheet_en.tex
\RequirePackage{fix-cm}
\documentclass[10pt,a4paper]{article}

\usepackage{amsmath,amssymb}
\usepackage{fontspec}
\setmainfont{cmunrm.otf}[BoldFont=cmunbx.otf, ItalicFont=cmunti.otf, BoldItalicFont=cmunbi.otf]
\setsansfont{cmunss.otf}[BoldFont=cmunsx.otf, ItalicFont=cmunsi.otf, BoldItalicFont=cmunso.otf]
\setmonofont{cmuntt.otf}[BoldFont=cmuntb.otf, ItalicFont=cmunit.otf, BoldItalicFont=cmuntx.otf]
\usepackage{polyglossia}
\setdefaultlanguage{english}

\usepackage[a4paper,left=8mm,right=8mm,top=7mm,bottom=11mm,footskip=5.5mm]{geometry}
\usepackage{xcolor}
\usepackage{multicol}
\usepackage{enumitem}
\usepackage{ragged2e}
\usepackage{fancyhdr}
\usepackage[colorlinks=true,linkcolor=hseblue,urlcolor=hseblue]{hyperref}

\definecolor{hseblue}{RGB}{16,45,105}
\definecolor{hseaccent}{RGB}{198,40,40}
\definecolor{hsegray}{RGB}{108,114,126}
\definecolor{hselight}{RGB}{240,243,248}

\hypersetup{pdftitle={Seminar 06 cheat sheet: Second-Order Methods, Part Two}, pdfauthor={Evgeny Baulin}}

\pagestyle{fancy}
\fancyhf{}
\renewcommand{\headrulewidth}{0pt}
\fancyfoot[L]{\scriptsize\color{hsegray}\textcopyright\ 2026 Evgeny Baulin, \href{https://t.me/tarakan_tuc}{t.me/tarakan\_tuc}}
\fancyfoot[R]{\scriptsize\color{hsegray}MIT License}

\setlength{\parindent}{0pt}
\setlength{\parskip}{0pt}
\setlength{\columnsep}{4.5mm}
\setlength{\columnseprule}{0.3pt}
\renewcommand{\columnseprulecolor}{\color{hselight}}
\setlength{\emergencystretch}{2em}
\setlength{\RaggedRightRightskip}{0pt plus 4em}
\binoppenalty=2000
\relpenalty=500
\setlength{\abovedisplayskip}{2pt}
\setlength{\belowdisplayskip}{2pt}
\setlist[itemize]{leftmargin=7pt, labelsep=3pt, itemsep=1.2pt, topsep=1pt, parsep=0pt, partopsep=0pt,
  label={\color{hseaccent}\scriptsize\textbullet}}

\makeatletter
\newcommand{\cssec}[1]{\par\addvspace{3.5pt}\noindent
  \colorbox{hselight}{\parbox{\dimexpr\linewidth-2\fboxsep}{\RaggedRight\hyphenpenalty=10000\color{hseblue}\bfseries\strut #1}}%
  \par\nobreak\vspace{1pt}\@afterheading}
\makeatother
\clubpenalty=10000
\widowpenalty=10000

\begin{document}
\thispagestyle{fancy}

\noindent
\begin{minipage}[b]{0.6\textwidth}
  \RaggedRight\hyphenpenalty=10000
  {\color{hseaccent}\bfseries\small Seminar 06 cheat sheet}\\[1pt]
  {\color{hseblue}\bfseries\Large Second-Order\\ Methods, Part Two}
\end{minipage}\hfill
\begin{minipage}[b]{0.39\textwidth}
  \raggedleft\footnotesize
  {\bfseries Optimization Methods}\\
  HSE University, Faculty of Computer Science, 2026--2027\\
  Course lecturer: Andrey Ignatov\\
  Author: Evgeny Baulin, Telegram \href{https://t.me/tarakan_tuc}{@tarakan\_tuc}
\end{minipage}

\vspace{2pt}
{\color{hseblue}\rule{\textwidth}{0.8pt}}
\vspace{-4pt}

\fontsize{9}{11.4}\selectfont
\begin{multicols}{3}
\RaggedRight


% ---- gn
\cssec{Gauss–Newton}
\begin{itemize}
  \item $f=\tfrac{1}{2}\|r\|^2$: $\nabla f=J^\top r$, $\nabla^2 f=J^\top J+\sum r_i\nabla^2 r_i$.
  \item GN: $J^\top Jp=-J^\top r$. Needs full column rank of $J$.
  \item Small residuals near $x^*$: quadratic convergence; big residuals: linear.
\end{itemize}

% ---- lm
\cssec{Levenberg–Marquardt}
\begin{itemize}
  \item $(J^\top J+\lambda I)p=-J^\top r$; $\rho=\frac{f-f^+}{m(0)-m(p)}$.
  \item $\rho\ge3/4$: accept, $\lambda\mathrel{/}=10$. $\rho<1/4$: reject, $\lambda\mathrel{*}=10$.
  \item Small $\lambda$: Gauss–Newton. Large $\lambda$: short gradient step.
\end{itemize}

% ---- tr
\cssec{Trust region}
\begin{itemize}
  \item $\min m(p)$ over $\|p\|\le\Delta$: $(B+\lambda I)p=-g$ with $\lambda\ge0$.
  \item Dogleg: broken line $0\to p_C\to p_N$; farthest point in the ball.
  \item $\rho>3/4$ on boundary: $\Delta\leftarrow 2\Delta$. $\rho<1/4$: $\Delta\leftarrow\Delta/4$.
\end{itemize}

% ---- lbfgs
\cssec{L-BFGS}
\begin{itemize}
  \item Store $m$ pairs $(s_i,y_i)$; $\gamma_k=s_{k-1}^\top y_{k-1}/(y_{k-1}^\top y_{k-1})$.
  \item Two-loop: back-sweep $a_i,q$; scale $\gamma_k q$; forward sweep $b_i,r$.
  \item One direction: $4mn$ work, $2mn$ memory, no $n\times n$ matrix.
\end{itemize}

\end{multicols}
\end{document}
